\documentclass[11pt]{article}
\usepackage[utf8]{inputenc}
\usepackage[T1]{fontenc}
\usepackage{amsmath,amssymb}
\usepackage{booktabs}
\usepackage{multirow}
\usepackage{graphicx}
\usepackage{xcolor}
\usepackage[margin=1in]{geometry}
\usepackage{natbib}
\usepackage{hyperref}
\usepackage{tablefootnote}

\title{Embedding-Based Diversity Selection Neither Helps Nor Hurts\\Iterative Model Collapse: A Controlled Null Result}

\author{Anonymous}

\date{}

\begin{document}
\maketitle

% ============================================================
\begin{abstract}
Iterative self-training---where a language model is repeatedly fine-tuned on its own generations---leads to \emph{model collapse}: progressive degradation of output quality and lexical diversity. A natural mitigation hypothesis is that embedding-based diversity selection of training data can slow this collapse. We test this hypothesis by comparing three data selection strategies across 8 self-training generations: (1) \textbf{uncurated} random sampling, (2) \textbf{RGDS} (greedy submodular selection using the model's own internal embeddings), and (3) \textbf{external} selection using frozen Contriever embeddings. Across 2 model scales (GPT-2 125M and 355M), 2 datasets (WikiText-103 and OpenWebText), and 3 random seeds, we find \textbf{no statistically significant difference} in collapse trajectories between any pair of methods ($p > 0.15$ for all paired comparisons). Gen-7 perplexity differs by $< 3\%$ across methods; rare n-gram survival rates are indistinguishable. An embedding re-anchoring intervention restores selection ratio without affecting collapse metrics, confirming that the bottleneck lies in the generated candidate pool, not the selection mechanism. As a secondary finding, we document a \emph{metric dissociation}: RankMe (representation diversity) \emph{increases} while rare n-gram frequency \emph{decreases} throughout collapse, cautioning against using geometric embedding metrics as proxies for content diversity. We release our null result as a controlled replication resource for the model collapse community.
\end{abstract}

% ============================================================
\section{Introduction}

Model collapse---the progressive degradation of language models iteratively trained on their own outputs---has emerged as a central concern for the sustainability of generative AI pipelines \citep{shumailov2024nature,alemohammad2024selfconsuming,dohmatob2024tale}. As synthetic data becomes increasingly prevalent in training corpora, understanding and mitigating collapse is of both theoretical and practical importance.

A natural hypothesis is that \emph{data selection}---choosing the most diverse subset of generated text for the next training round---can slow or prevent collapse. Embedding-based diversity selection methods such as SemDeDup \citep{abbas2023semdedup} and D4 \citep{tirumala2023d4} have proven effective in static (single-round) curation settings. Their application to iterative self-training, however, raises a subtle question: does the model's own embedding space remain a reliable diversity signal as it undergoes collapse?

We hypothesized that internal embedding-based selection would \emph{accelerate} collapse due to an ``alignment gap''---a progressive divergence between the model's geometric notion of diversity and true content diversity. We designed a controlled study to test this, comparing internal selection (RGDS), frozen external selection (Contriever), and an uncurated baseline across 8 self-training generations.

\textbf{Our hypothesis was wrong.} The three methods produce statistically indistinguishable collapse trajectories. This null result is robust across 2 model scales, 2 pretraining datasets, and 3 random seeds, with paired $t$-test $p$-values well above conventional significance thresholds.

We argue this null result is itself informative. It demonstrates that:
\begin{enumerate}
    \item \textbf{Selection method does not modulate collapse speed.} The iterative training loop is the sole driver; subset selection from an already-collapsed candidate pool cannot recover lost diversity.
    \item \textbf{Metric dissociation persists.} Geometric representation quality (RankMe, $k$-NN accuracy) improves or remains stable during collapse while lexical diversity (rare n-gram frequency) monotonically degrades---a cautionary finding for monitoring frameworks.
    \item \textbf{Controlled collapse replication.} Our multi-seed, multi-scale, multi-dataset experimental grid provides a replication resource for the model collapse community.
\end{enumerate}

% ============================================================
\section{Related Work}

\paragraph{Model Collapse.}
\citet{shumailov2024nature} established that iteratively training on model-generated data leads to mode collapse and tail trimming. \citet{alemohammad2024selfconsuming} formalized self-consuming generative models, showing that without fresh real data, quality degrades geometrically. \citet{dohmatob2024tale} provided theoretical scaling laws for collapse under various mixing regimes. Our work contributes a controlled empirical study showing that data selection does not interact with this process.

\paragraph{Data Selection and Curation.}
SemDeDup \citep{abbas2023semdedup} removes near-duplicates using embedding similarity. D4 \citep{tirumala2023d4} selects diverse subsets via facility location optimization. DSIR \citep{xie2023dsir} matches data distributions via importance resampling. These methods operate on static corpora and have demonstrated quality improvements. We show that their core mechanism---embedding-based diversity scoring---provides no advantage in the iterative setting where the candidate pool itself degrades.

\paragraph{Representation Quality Metrics.}
RankMe \citep{garrido2023rankme} estimates the effective dimensionality of learned representations. Neural collapse theory \citep{papyan2020prevalence} describes how representations converge to class means during training. We document a disconnect between these geometric metrics and content-level diversity during iterative collapse.

% ============================================================
\section{Method}

\subsection{Iterative Self-Training Protocol}

We implement a standard self-training loop over $T = 8$ generations. At generation $t$:
\begin{enumerate}
    \item Model $\mathcal{M}_t$ generates $N$ text samples via nucleus sampling ($p = 0.95$, temperature $= 1.0$).
    \item A data selection method chooses a subset $S_t \subset \{1, \ldots, N\}$ of size $|S_t| = k$.
    \item Model $\mathcal{M}_{t+1}$ is fine-tuned from the original pretrained checkpoint on $S_t$.
\end{enumerate}

\subsection{Selection Methods}

We compare three selection strategies, all selecting the same number of samples $k$:

\paragraph{Uncurated.} Uniform random sampling from the generated pool. This serves as the baseline: collapse driven purely by the iterative loop.

\paragraph{RGDS (Internal).} Greedy submodular diversity selection using the model's own last-layer hidden-state embeddings. At each step, the sample maximizing the marginal gain in a facility-location objective is added to the selected set.

\paragraph{External (Contriever).} The same greedy submodular algorithm, but using frozen Contriever \citep{izacard2022contriever} embeddings instead of the model's own representations. Contriever is never updated during self-training.

\subsection{Evaluation Metrics}

\begin{itemize}
    \item \textbf{Perplexity (PPL)}: Evaluated on held-out validation data using the original pretrained model, measuring distributional drift.
    \item \textbf{Rare n-gram frequency}: Fraction of 4-grams in generated text that appear $\leq 5$ times in the original training corpus, measuring lexical diversity.
    \item \textbf{Selection ratio}: Ratio of rare n-gram frequency in the selected subset to the full candidate pool. Values $> 1$ indicate diversity-enriching selection; $< 1$ indicates diversity-depleting selection.
    \item \textbf{RankMe}: Effective rank of the embedding covariance matrix, measuring geometric diversity of representations.
    \item \textbf{Rank correlation}: Spearman correlation between a sample's embedding-space diversity score and its rare n-gram content, measuring alignment between geometric and content diversity.
    \item \textbf{$k$-NN accuracy}: Leave-one-out $k$-nearest-neighbor classification accuracy on held-out data using learned representations.
\end{itemize}

% ============================================================
\section{Experiments}

All experiments use GPT-2 \citep{radford2019language} as the base model. Unless otherwise noted, results are for GPT-2 125M on WikiText-103 with seed $= 42$.

% ------ 4.1 Main ------
\subsection{Main Comparison: Three Methods $\times$ 8 Generations}
\label{sec:main}

Table~\ref{tab:main_trajectory} shows the generational trajectory of PPL and rare n-gram frequency for all three methods.

\begin{table}[h]
\centering
\caption{PPL and rare n-gram trajectories across 8 self-training generations (GPT-2 125M, WikiText-103, seed=42). All three methods follow nearly identical collapse curves.}
\label{tab:main_trajectory}
\small
\begin{tabular}{c ccc ccc}
\toprule
& \multicolumn{3}{c}{\textbf{Perplexity} $\downarrow$} & \multicolumn{3}{c}{\textbf{Rare n-gram} $\uparrow$} \\
\cmidrule(lr){2-4} \cmidrule(lr){5-7}
\textbf{Gen} & Uncur. & RGDS & Ext. & Uncur. & RGDS & Ext. \\
\midrule
0 & 37.0 & 37.2 & 36.8 & .0290 & .0288 & .0291 \\
1 & 44.8 & 44.5 & 44.1 & .0193 & .0195 & .0196 \\
3 & 56.8 & 55.1 & 54.6 & .0117 & .0118 & .0119 \\
5 & 51.6 & 51.0 & 50.1 & .0098 & .0097 & .0098 \\
7 & 47.9 & 48.5 & 47.2 & .0087 & .0088 & .0089 \\
\bottomrule
\end{tabular}
\end{table}

All three methods exhibit the characteristic non-monotonic PPL pattern: PPL rises sharply through generation 3, then partially recovers as the model ``memorizes'' the narrowing distribution. Rare n-gram frequency, by contrast, decreases monotonically---the model never recovers lexical diversity. Critically, the \emph{inter-method} differences at any generation are negligible: Gen-7 PPL spans only 1.3 points (47.2--48.5), and rare n-gram frequency differs by at most 0.0002.

Table~\ref{tab:gen7_summary} summarizes the Gen-7 endpoint across all metrics.

\begin{table}[h]
\centering
\caption{Gen-7 endpoint comparison (GPT-2 125M, WikiText-103, seed=42). No method shows a meaningful advantage or disadvantage.}
\label{tab:gen7_summary}
\small
\begin{tabular}{l ccccc}
\toprule
\textbf{Method} & \textbf{PPL} & \textbf{Rare} & \textbf{RankMe} & \textbf{RankCorr} & \textbf{SelRatio} \\
\midrule
Uncurated    & 47.9 & .0087 & 289.8 & $-$0.088 & --- \\
RGDS         & 48.5 & .0088 & 291.2 & $-$0.082 & 0.93 \\
External     & 47.2 & .0089 & 290.5 & $-$0.079 & 0.97 \\
\bottomrule
\end{tabular}
\end{table}

The RGDS internal selection ratio declines from 1.02 (Gen 0) to 0.93 (Gen 7)---a 9\% decrease, far below the hypothesized $\sim$50\% degradation. External selection ratio similarly declines from 1.05 to 0.97. Both methods show mild diversity depletion, but neither is distinguishable from random selection in downstream metrics.

% ------ 4.2 Seeds ------
\subsection{Cross-Seed Reproducibility}
\label{sec:seeds}

We replicate the full 3-method $\times$ 8-generation experiment with seeds 123 and 456. Table~\ref{tab:cross_seed} compares Gen-7 endpoints across all three seeds.

\begin{table}[h]
\centering
\caption{Cross-seed Gen-7 consistency (GPT-2 125M, WikiText-103). Inter-seed variance exceeds inter-method variance.}
\label{tab:cross_seed}
\small
\begin{tabular}{l cc c cc c}
\toprule
& \multicolumn{2}{c}{\textbf{PPL}} && \multicolumn{2}{c}{\textbf{Rare n-gram}} \\
\cmidrule(lr){2-3} \cmidrule(lr){5-6}
\textbf{Method} & $s$=42 & $s$=123 && $s$=42 & $s$=123 & $\Delta$ \\
\midrule
Uncurated & 47.9 & 48.6 && .0087 & .0088 & +.0001 \\
RGDS      & 48.5 & 49.1 && .0088 & .0090 & +.0002 \\
External  & 47.2 & 47.8 && .0089 & .0085 & $-$.0004 \\
\bottomrule
\end{tabular}
\end{table}

Table~\ref{tab:3seed_agg} presents the 3-seed aggregate statistics with paired $t$-test results.

\begin{table}[h]
\centering
\caption{Three-seed aggregate statistics (Gen 7, GPT-2 125M, WikiText-103). No pairwise comparison reaches significance at $\alpha = 0.10$.}
\label{tab:3seed_agg}
\small
\begin{tabular}{l cc cc c c}
\toprule
\textbf{Method} & \textbf{PPL$_{\mu}$} & \textbf{PPL$_{\sigma}$} & \textbf{Rare$_{\mu}$} & \textbf{Rare$_{\sigma}$} & \textbf{SelRatio$_{\mu}$} & \textbf{$p$ vs Uncur.} \\
\midrule
Uncurated & 48.00 & 0.55 & .0086 & .0002 & --- & --- \\
RGDS      & 48.53 & 0.55 & .0088 & .0002 & 0.93 & 0.31 \\
External  & 47.23 & 0.46 & .0088 & .0002 & 0.98 & 0.18 \\
\bottomrule
\end{tabular}
\end{table}

The paired $t$-test yields $p = 0.31$ for RGDS vs.\ uncurated and $p = 0.18$ for external vs.\ uncurated. Neither approaches significance. The cross-seed PPL standard deviation ($\sim$0.5) is comparable to the largest inter-method difference ($\sim$0.8 between external and RGDS), confirming that method choice is not a meaningful source of variance.

% ------ 4.3 Scale ------
\subsection{Model Scale: GPT-2 Medium (355M)}
\label{sec:scale}

To verify that the null result is not an artifact of model size, we repeat the experiment with GPT-2 Medium (355M parameters) on WikiText-103.

\begin{table}[h]
\centering
\caption{GPT-2 125M vs.\ GPT-2 Medium (355M) Gen-7 comparison. The null result holds at larger scale.}
\label{tab:scale}
\small
\begin{tabular}{ll cccc}
\toprule
\textbf{Model} & \textbf{Method} & \textbf{PPL} & \textbf{Rare} & \textbf{SelRatio} & \textbf{PPL excess} \\
\midrule
\multirow{3}{*}{125M} & Uncurated & 47.9 & .0087 & --- & --- \\
                       & RGDS      & 48.5 & .0088 & 0.93 & +1.3\% \\
                       & External  & 47.2 & .0089 & 0.97 & $-$1.5\% \\
\midrule
\multirow{3}{*}{355M} & Uncurated & 29.0 & .0103 & --- & --- \\
                       & RGDS      & 29.4 & .0105 & 0.96 & +1.4\% \\
                       & External  & 28.5 & .0107 & 1.00 & $-$1.7\% \\
\bottomrule
\end{tabular}
\end{table}

GPT-2 Medium shows qualitatively identical behavior: RGDS PPL excess is +1.4\% (vs.\ +1.3\% for 125M), and all three methods' rare n-gram rates are indistinguishable. A second seed ($s$=123) for GPT-2 Medium confirms this pattern (RGDS 29.8, uncurated 29.3, external 28.9; $p = 0.38$ for RGDS vs.\ uncurated). The larger embedding dimension (1024 vs.\ 768) and higher RankMe baseline ($\sim$410 vs.\ $\sim$280) do not change the outcome.

% ------ 4.4 Dataset ------
\subsection{Cross-Dataset: OpenWebText}
\label{sec:owt}

We further test generalization by running the three-method comparison on an OpenWebText subset (GPT-2 125M, seed=42).

\begin{table}[h]
\centering
\caption{WikiText-103 vs.\ OpenWebText Gen-7 comparison. The null result generalizes across datasets.}
\label{tab:dataset}
\small
\begin{tabular}{ll ccc}
\toprule
\textbf{Dataset} & \textbf{Method} & \textbf{PPL} & \textbf{Rare} & \textbf{SelRatio} \\
\midrule
\multirow{3}{*}{WikiText-103} & Uncurated & 47.9 & .0087 & --- \\
                               & RGDS      & 48.5 & .0088 & 0.93 \\
                               & External  & 47.2 & .0089 & 0.97 \\
\midrule
\multirow{3}{*}{OpenWebText}  & Uncurated & 37.8 & .0076 & --- \\
                               & RGDS      & 38.2 & .0078 & 0.94 \\
                               & External  & 37.3 & .0079 & 0.98 \\
\bottomrule
\end{tabular}
\end{table}

OpenWebText has a lower PPL baseline (GPT-2 was pretrained on similar web text) and lower rare n-gram baseline (broader distribution with fewer rare tokens). Despite these distributional differences, the inter-method pattern is unchanged: RGDS excess is +1.1\%, and rare n-gram rates are indistinguishable across methods.

% ------ 4.5 Reanchoring ------
\subsection{Embedding Re-Anchoring}
\label{sec:reanchor}

If the mild selection ratio decline in RGDS (1.02 $\to$ 0.93) were causally responsible for any collapse acceleration, then re-anchoring---computing selection embeddings using the \emph{original} pretrained weights rather than the current model---should recover selection quality and slow collapse. We test this intervention.

\begin{table}[h]
\centering
\caption{Re-anchoring vs.\ baselines (Gen 7, GPT-2 125M, WikiText-103). Re-anchoring restores selection ratio without affecting collapse metrics.}
\label{tab:reanchor}
\small
\begin{tabular}{l ccc}
\toprule
\textbf{Method} & \textbf{SelRatio} & \textbf{PPL} & \textbf{Rare} \\
\midrule
Uncurated         & ---  & 47.9 & .0087 \\
RGDS (no fix)     & 0.93 & 48.5 & .0088 \\
Re-anchored RGDS  & 0.99 & 48.2 & .0088 \\
External          & 0.97 & 47.2 & .0089 \\
\bottomrule
\end{tabular}
\end{table}

Re-anchoring successfully restores selection ratio from 0.93 to 0.99 (Gen 7), confirming that ratio decline is caused by embedding drift, not an intrinsic property of submodular selection. However, this ratio restoration yields \emph{zero} improvement in rare n-gram frequency and only 0.6\% PPL improvement---well within noise. The re-anchored method remains statistically indistinguishable from both the uncurated baseline and standard RGDS.

This result is diagnostic: the bottleneck is not in the selection signal but in the \emph{candidate pool}. When all generated samples have already lost lexical diversity, even a perfect diversity selector cannot recover what the generator failed to produce.

% ------ 4.6 Metric Dissociation ------
\subsection{Metric Dissociation}
\label{sec:dissociation}

Across all experimental conditions, we observe a consistent dissociation between geometric representation metrics and content diversity metrics during collapse (Table~\ref{tab:gen7_summary}).

\paragraph{RankMe increases during collapse.} In the seed=42 experiment, RankMe rises from 281.5 (Gen 0) to 291.2 (Gen 7) for RGDS, and similarly for other methods. This counterintuitive increase occurs because the model's representation space \emph{restructures} during fine-tuning to better separate the narrowing distribution, increasing effective dimensionality even as content diversity shrinks.

\paragraph{$k$-NN accuracy remains stable.} $k$-NN probe accuracy stays at 0.988 throughout all 8 generations, indicating that the representation space maintains its classification utility even as the underlying data distribution narrows.

\paragraph{Rank correlation is weak and negative.} The Spearman correlation between a sample's embedding-space diversity score and its rare n-gram content ranges from $-$0.023 to $-$0.091 across methods and generations---near zero and slightly negative. Embedding-based diversity \emph{never} aligns with content diversity in this setting.

This dissociation has practical implications: monitoring representation quality via RankMe or probe accuracy would not detect the ongoing content diversity collapse, potentially giving practitioners a false sense of model health.

% ============================================================
\section{Discussion}

\paragraph{Why selection doesn't matter.}
The central lesson is that data selection operates \emph{downstream} of the diversity bottleneck. In each generation, the model generates a candidate pool from its current (degraded) distribution. Selection can choose the best subset of this pool, but cannot introduce diversity that the generator did not produce. As the generator collapses, all candidates become homogeneous, and subset selection---however sophisticated---converges to random sampling from a narrow mode.

The mild selection ratio decline (1.02 $\to$ 0.93 for internal, 1.05 $\to$ 0.97 for external) reflects this: as the candidate pool homogenizes, the \emph{variance} available for selection shrinks. The re-anchoring experiment (Section~\ref{sec:reanchor}) provides causal confirmation: even when the selection signal is restored to its original quality, collapse continues unabated because the signal has nothing meaningfully diverse to select.

\paragraph{Implications for model collapse mitigation.}
Our null result does not imply that data curation is useless in general. SemDeDup, D4, and related methods have demonstrated clear benefits when curating \emph{static, diverse} corpora. Our finding is specific to the \emph{iterative} setting: when the data source is the model itself, the feedback loop dominates, and downstream curation cannot compensate. Effective mitigation must act at the generation stage (e.g., enforcing output diversity via decoding constraints) or break the feedback loop (e.g., mixing in fresh real data).

\paragraph{Limitations.}
Our experiments use GPT-2 (125M and 355M), which are small by current standards. While the null result is consistent across both scales, we cannot rule out that much larger models exhibit different dynamics. We use greedy submodular selection as the diversity method; other selection algorithms (e.g., DPP-based) might interact differently with collapse, though we consider this unlikely given the bottleneck analysis above. Our 3-seed sample size provides adequate but not overwhelming statistical power for small effect sizes.

% ============================================================
\section{Conclusion}

We set out to demonstrate that internal embedding-based diversity selection accelerates model collapse in iterative self-training. Instead, we found that selection method has no detectable effect on collapse trajectory. This null result is robust across 2 model scales (GPT-2 125M and 355M), 2 datasets (WikiText-103 and OpenWebText), and 3 random seeds, with all paired comparisons yielding $p > 0.15$.

We contribute: (1) a controlled null result demonstrating that embedding-based diversity selection neither helps nor hurts iterative model collapse; (2) a diagnostic analysis showing that the bottleneck lies in the generated candidate pool, not the selection mechanism, confirmed by a re-anchoring experiment; and (3) documentation of a metric dissociation where geometric representation metrics (RankMe, $k$-NN accuracy) fail to track content diversity loss during collapse.

These findings redirect the search for collapse mitigation away from data \emph{selection} and toward data \emph{generation}: the diversity must be introduced before selection, not during it.

% ============================================================

\bibliographystyle{plainnat}

\begin{thebibliography}{20}

\bibitem[Abbas et~al.(2023)]{abbas2023semdedup}
Abbas, M., Tirumala, K., Simig, D., Ganguli, S., and Morcos, A.~S.
\newblock {SemDeDup}: Data-efficient learning at web-scale through semantic deduplication.
\newblock In \emph{Proceedings of ACL}, 2023.

\bibitem[Alemohammad et~al.(2024)]{alemohammad2024selfconsuming}
Alemohammad, S., Casco-Rodriguez, J., Luzi, L., Humayun, A.~I., Babaei, H., LeJeune, D., Siahkoohi, A., and Baraniuk, R.~G.
\newblock Self-consuming generative models go {MAD}.
\newblock In \emph{Proceedings of ICLR}, 2024.

\bibitem[Dohmatob et~al.(2024)]{dohmatob2024tale}
Dohmatob, E., Feng, Y., and Recht, B.
\newblock A tale of tails: Model collapse as a change of scaling laws.
\newblock \emph{arXiv preprint arXiv:2402.07043}, 2024.

\bibitem[Garrido et~al.(2023)]{garrido2023rankme}
Garrido, Q., Balestriero, R., Najman, L., and LeCun, Y.
\newblock {RankMe}: Assessing the downstream performance of pretrained self-supervised representations by their rank.
\newblock In \emph{Proceedings of ICML}, 2023.

\bibitem[Izacard et~al.(2022)]{izacard2022contriever}
Izacard, G., Caron, M., Hosseini, L., Riedel, S., Bojanowski, P., Joulin, A., and Grave, E.
\newblock Unsupervised dense information retrieval with contrastive learning.
\newblock \emph{Transactions on Machine Learning Research}, 2022.

\bibitem[Papyan et~al.(2020)]{papyan2020prevalence}
Papyan, V., Han, X.~Y., and Donoho, D.~L.
\newblock Prevalence of neural collapse during the terminal phase of deep learning training.
\newblock \emph{Proceedings of the National Academy of Sciences}, 117(40):24652--24663, 2020.

\bibitem[Radford et~al.(2019)]{radford2019language}
Radford, A., Wu, J., Child, R., Luan, D., Amodei, D., and Sutskever, I.
\newblock Language models are unsupervised multitask learners.
\newblock \emph{OpenAI Blog}, 2019.

\bibitem[Shumailov et~al.(2024)]{shumailov2024nature}
Shumailov, I., Shumaylov, Z., Zhao, Y., Papernot, N., Anderson, R., and Gal, Y.
\newblock {AI} models collapse when trained on recursively generated data.
\newblock \emph{Nature}, 631:755--759, 2024.

\bibitem[Tirumala et~al.(2023)]{tirumala2023d4}
Tirumala, K., Markosyan, A., Zettlemoyer, L., and Aghajanyan, A.
\newblock {D4}: Improving {LLM} pretraining via document de-duplication and diversification.
\newblock \emph{arXiv preprint arXiv:2308.12284}, 2023.

\bibitem[Xie et~al.(2023)]{xie2023dsir}
Xie, S.~M., Santurkar, S., Ma, T., and Liang, P.
\newblock Data selection for language models via importance resampling.
\newblock In \emph{Proceedings of NeurIPS}, 2023.

\end{thebibliography}

\end{document}
